\documentclass{article}
\usepackage{amsmath, amsthm, amssymb}
\usepackage{geometry}
\geometry{margin=1in}

\newtheorem{lemma}{Lemma}[section]
\newtheorem{theorem}[lemma]{Theorem}

\title{}
\date{}

\begin{document}

\section{Large $\varepsilon$-light vertex subsets}

\subsection*{Problem}

For a graph $G = (V, E)$, let $G_S = (V, E(S,S))$ denote the graph with the same vertex set,
but only the edges between vertices in $S$. Let $L$ be the Laplacian matrix of $G$ and let $L_S$
be the Laplacian of $G_S$. I say that a set of vertices $S$ is $\epsilon$-light if the matrix
$\epsilon L - L_S$ is positive semidefinite. Does there exist a constant $c > 0$ so that for every
graph $G$ and every $\epsilon$ between $0$ and $1$, $V$ contains an $\epsilon$-light subset $S$ of size
at least $c\epsilon|V|$?

\subsection*{Solution}

We prove the claim with an explicit constant $c = 1/256$.

Throughout we write $n = |V|$. Some degenerate cases are immediate and will be set aside.
If $n = 0$ the assertion is vacuous. If $\varepsilon = 0$ the empty set is $0$-light and has the required
size. If the graph has no edges (so that its Laplacian is the zero matrix), then $L_S = 0$ for
every $S$ and we may simply take all vertices. Thus in the main part of the proof we assume
$n \ge 1$, $\varepsilon \in (0,1]$, and that the Laplacian has positive rank. All Loewner inequalities
and traces below are taken on the subspace $\operatorname{range}(L) = (\ker L)^\perp$, and $I$ denotes the
identity on that space. The exceptional cases are revisited in Step~6.

\paragraph{Step 1: Normalization on the Laplacian range.}
Let $\ker L$ be the space of vectors that are constant on each connected component of $G$.
Let $L^\dagger$ be the Moore--Penrose pseudoinverse, and define
\[
  L^{-1/2} := (L^\dagger)^{1/2}.
\]
Then $L^{-1/2}$ acts as the inverse square root on $\operatorname{range}(L) = (\ker L)^\perp$ and as $0$ on $\ker L$.

For an edge $e = \{u,v\}$ define the rank-one edge Laplacian
\[
  L_e := (e_u - e_v)(e_u - e_v)^\top,
\]
so that $L = \sum_{e \in E} L_e$. All sums over edges below are taken with multiplicity if the
graph has parallel edges. Define
\[
  A_e := L^{-1/2} L_e L^{-1/2}.
\]
Each $A_e$ is positive semidefinite on $\operatorname{range}(L)$, whose dimension is $d := \operatorname{rank}(L) \le n$.
Moreover, on $\operatorname{range}(L)$ we have
\begin{equation}
  \sum_{e \in E} A_e = L^{-1/2}\!\left(\sum_{e \in E} L_e\right)\!L^{-1/2} = L^{-1/2} L L^{-1/2} = I. \label{eq:sum1}
\end{equation}
Also, for any $S \subseteq V$,
\begin{equation}
  L^{-1/2} L_S L^{-1/2} = \sum_{e \in E(S,S)} A_e \quad \text{on } \operatorname{range}(L). \label{eq:LS}
\end{equation}
Therefore, it suffices to find $S$ such that on $\operatorname{range}(L)$,
\begin{equation}
  \sum_{e \in E(S,S)} A_e \preceq \varepsilon I, \label{eq:goal}
\end{equation}
because then \eqref{eq:LS} implies $L^{-1/2} L_S L^{-1/2} \preceq \varepsilon I$, i.e.
\[
  x^\top L_S x \le \varepsilon\, x^\top L x \quad \text{for all } x \perp \ker L.
\]
If $P = L^{-1/2}L^{1/2} = L^{1/2}L^{-1/2}$ denotes the orthogonal projection onto $\operatorname{range}(L)$, the
displayed Loewner inequality is equivalent to $x^T L_S x \le \varepsilon x^T L x$ for every
$x \in \operatorname{range}(L)$ by taking $z = L^{1/2}x$ (where $L^{1/2}$ denotes the positive square root of $L$,
acting as zero on $\ker L$) in the quadratic form. Vectors in $\ker L$ are constant on each
connected component and are therefore also annihilated by $L_S$; by symmetry no mixed terms
occur between $\operatorname{range}(L)$ and $\ker L$. Hence the inequality holds for all $x \in \mathbb{R}^V$, which
is exactly $\varepsilon L - L_S \succeq 0$.

\paragraph{Step 2: A one-sided BSS barrier lemma.}
The following lemma is a one-sided variant of the barrier method introduced by Batson, Spielman
and Srivastava~\cite{BSS}; we give a complete proof for the reader's convenience. For a PSD
matrix $M \succeq 0$ on a $d$-dimensional space and a scalar $u > \lambda_{\max}(M)$, define
the potential
\[
  \Phi_u(M) := \operatorname{tr}\!\bigl((uI - M)^{-1}\bigr).
\]

\begin{lemma}[One-sided barrier]\label{lem:barrier}
Assume $M \prec uI$, let $u' > u$, and put $U := (u'I - M)^{-1}$. If $B \succeq 0$ satisfies
\begin{equation}
  \operatorname{tr}(BU) + \frac{\operatorname{tr}(BU^2)}{\Phi_u(M) - \Phi_{u'}(M)} \le 1, \label{eq:barcond}
\end{equation}
then $M + B \prec u'I$ and $\Phi_{u'}(M + B) \le \Phi_u(M)$.
\end{lemma}

\begin{proof}
Let $K := B^{1/2} U B^{1/2} \succeq 0$. The hypothesis \eqref{eq:barcond} implies $\operatorname{tr}(K) < 1$:
the second summand there is non-negative, and if it were zero then the positive semidefinite
matrix $B^{1/2} U^2 B^{1/2}$ would have trace zero and hence vanish; since $U$ is invertible on
our space this forces $B = 0$. Consequently every eigenvalue of $K$ is $< 1$, so in particular
$\|K\| < 1$ and $(I - K)$ is invertible. By the Sherman--Morrison--Woodbury identity (which
can be verified by multiplying the two sides),
\[
  (u'I - M - B)^{-1} = U + U B^{1/2}(I - K)^{-1} B^{1/2} U,
\]
so $u'I - M - B \succ 0$, i.e.\ $M + B \prec u'I$.

Taking traces, using cyclicity of the trace and the elementary fact that $\operatorname{tr}(XC) \le \operatorname{tr}(YC)$
whenever $0 \preceq X \preceq Y$ and $C \succeq 0$, together with $(I-K)^{-1} \preceq (1 - \operatorname{tr} K)^{-1} I$
(valid for PSD $K$ with $\operatorname{tr} K < 1$), we obtain
\[
  \Phi_{u'}(M+B) \le \Phi_{u'}(M) + \frac{\operatorname{tr}(BU^2)}{1 - \operatorname{tr}(BU)}.
\]
A short rearrangement shows that \eqref{eq:barcond} is equivalent to the bound that the
right-hand side is at most $\Phi_u(M)$. This yields $\Phi_{u'}(M+B) \le \Phi_u(M)$.
\end{proof}

We will also use the following inequality: if $u' = u + \delta$, then
\begin{equation}
  \Phi_u(M) - \Phi_{u'}(M) \ge \delta\, \operatorname{tr}\!\bigl((u'I - M)^{-2}\bigr) = \delta\, \operatorname{tr}(U^2). \label{eq:drop}
\end{equation}
Indeed, diagonalizing $M$ with eigenvalues $\lambda_j$ gives
\[
  \Phi_u(M) - \Phi_{u'}(M) = \sum_j \frac{\delta}{(u - \lambda_j)(u' - \lambda_j)}
  \ge \sum_j \frac{\delta}{(u' - \lambda_j)^2}.
\]

\paragraph{Step 3: A partial coloring process.}
Fix $\varepsilon \in (0,1]$ and set
\begin{equation}
  r := \left\lceil \frac{16}{\varepsilon} \right\rceil, \quad
  u_0 := \frac{\varepsilon}{2}, \quad
  \delta := \frac{\varepsilon}{n}, \quad
  k := \left\lfloor \frac{n}{4} \right\rfloor. \label{eq:params}
\end{equation}
We will color $k$ vertices, one at a time, using $r$ colors.

At time $t$ ($0 \le t \le k$), let $T \subseteq V$ be the set of colored vertices, $|T| = t$, and
$\mathrm{col} : T \to \{1, \ldots, r\}$ the coloring. Define the PSD matrix (on $\operatorname{range}(L)$)
\begin{equation}
  M_t := \sum_{\substack{uv \in E \\ u,v \in T \\ \mathrm{col}(u) = \mathrm{col}(v)}} A_{uv}. \label{eq:Mt}
\end{equation}
Thus $M_t$ contains the contributions from edges whose endpoints are already colored and share
the same color.

Let $R := V \setminus T$ be the uncolored vertices, $m := |R| = n - t$. For $v \in R$ and
$\gamma \in \{1, \ldots, r\}$ define the prospective increment obtained by coloring $v$ with $\gamma$:
\begin{equation}
  B_v^\gamma := \sum_{\substack{u \in T \\ \mathrm{col}(u) = \gamma \\ uv \in E}} A_{uv}. \label{eq:Bvg}
\end{equation}
Then if we color $v$ with color $\gamma$, we have $M_{t+1} = M_t + B_v^\gamma$.

\paragraph{Step 4: Inductive barrier invariant.}
Let $u_t := u_0 + t\delta$. We maintain the invariant
\begin{equation}
  M_t \prec u_t I \quad \text{and} \quad \Phi_{u_t}(M_t) \le \Phi_{u_0}(0) = \frac{d}{u_0}. \label{eq:inv}
\end{equation}
This holds at $t = 0$ since $M_0 = 0$.

Assume it holds for some $t < k$. Set $u = u_t$, $u' = u_{t+1} = u_t + \delta$, and
\[
  U := (u'I - M_t)^{-1}.
\]
We claim there exists a choice of $(v, \gamma) \in R \times \{1, \ldots, r\}$ for which the barrier
condition \eqref{eq:barcond} holds with $M = M_t$ and $B = B_v^\gamma$.

Consider the average over a uniformly random pair $(v, \gamma)$:
\begin{equation}
  \frac{1}{mr} \sum_{v \in R} \sum_{\gamma=1}^r \left[
    \operatorname{tr}(B_v^\gamma U) + \frac{\operatorname{tr}(B_v^\gamma U^2)}{\Phi_u(M_t) - \Phi_{u'}(M_t)}
  \right]. \label{eq:avg}
\end{equation}
Observe that
\[
  \sum_{v \in R} \sum_{\gamma=1}^r B_v^\gamma
  = \sum_{\substack{uv \in E \\ u \in T,\, v \in R}} A_{uv} \preceq \sum_{e \in E} A_e = I
  \quad \text{on } \operatorname{range}(L),
\]
because the left-hand side is a sub-sum of the PSD matrices $\{A_e\}$ in \eqref{eq:sum1}. If
$X \preceq Y$ and $C \succeq 0$, then $\operatorname{tr}(XC) \le \operatorname{tr}(YC)$ because
$\operatorname{tr}(C^{1/2}(Y-X)C^{1/2}) \ge 0$. Applying this observation with $C = U$ and with $C = U^2$
(both positive semidefinite) we get
\[
  \sum_{v,\gamma} \operatorname{tr}(B_v^\gamma U) \le \operatorname{tr}(U), \qquad
  \sum_{v,\gamma} \operatorname{tr}(B_v^\gamma U^2) \le \operatorname{tr}(U^2).
\]
Therefore \eqref{eq:avg} is at most
\begin{equation}
  \frac{\operatorname{tr}(U)}{mr} + \frac{\operatorname{tr}(U^2)}{mr\bigl(\Phi_u(M_t) - \Phi_{u'}(M_t)\bigr)}. \label{eq:avg2}
\end{equation}
By the inductive hypothesis, $\operatorname{tr}(U) = \Phi_{u'}(M_t) \le \Phi_u(M_t) \le d/u_0$; the middle
inequality uses that, for fixed $M_t$, the function $s \mapsto \Phi_s(M_t)$ decreases as the
barrier level $s$ increases. By \eqref{eq:drop},
\[
  \Phi_u(M_t) - \Phi_{u'}(M_t) \ge \delta\, \operatorname{tr}(U^2),
\]
so (in the non-trivial case $d > 0$, where $\operatorname{tr}(U^2) > 0$) the second term in \eqref{eq:avg2}
is at most $1/(\delta m r)$. Hence the average \eqref{eq:avg} is at most
\begin{equation}
  \frac{d/u_0}{mr} + \frac{1}{\delta m r}. \label{eq:avg3}
\end{equation}
As long as $t < k = \lfloor n/4 \rfloor$, we have $m = n - t \ge 3n/4$ and $d \le n$. Using the
choices \eqref{eq:params}, we bound
\[
  \frac{d/u_0}{mr} \le \frac{n/(\varepsilon/2)}{(3n/4)\cdot(16/\varepsilon)} = \frac{1}{6}, \qquad
  \frac{1}{\delta m r} \le \frac{1}{(\varepsilon/n)\cdot(3n/4)\cdot(16/\varepsilon)} = \frac{1}{12},
\]
so the average \eqref{eq:avg3} is $< 1$. Therefore there exists at least one pair $(v, \gamma)$
for which \eqref{eq:barcond} holds. Applying Lemma~\ref{lem:barrier} yields
\[
  M_{t+1} \prec u_{t+1} I \quad \text{and} \quad \Phi_{u_{t+1}}(M_{t+1}) \le \Phi_{u_t}(M_t) \le \frac{d}{u_0}.
\]
Thus the invariant \eqref{eq:inv} propagates to $t+1$, completing the induction for
$t = 0, 1, \ldots, k$.

\paragraph{Step 5: Extracting a large $\varepsilon$-light set.}
After $k$ steps, the colored set $T$ (with $|T| = k$) is partitioned into $r$ color classes
$S_1, \ldots, S_r$. By definition of $M_k$,
\[
  M_k = \sum_{a=1}^r \sum_{uv \in E(S_a, S_a)} A_{uv} = \sum_{a=1}^r L^{-1/2} L_{S_a} L^{-1/2}
  \quad \text{on } \operatorname{range}(L).
\]
From the invariant, $M_k \preceq u_k I$ with $u_k = u_0 + k\delta \le \varepsilon/2 + \varepsilon/4 = 3\varepsilon/4$.
Since each summand is PSD, each is dominated by the sum. Let $S$ be the largest color class.
Then
\[
  L^{-1/2} L_S L^{-1/2} \preceq M_k \preceq \frac{3\varepsilon}{4}\,I \preceq \varepsilon I
  \quad \text{on } \operatorname{range}(L).
\]
As explained in Step~1, this implies $L_S \preceq \varepsilon L$, i.e.\ $\varepsilon L - L_S \succeq 0$, so $S$
is $\varepsilon$-light.

\paragraph{Step 6: Size lower bound.}
Among the $k$ colored vertices, the largest color class has size at least $k/r$. If $n \ge 4$,
then $k = \lfloor n/4 \rfloor \ge n/8$. Also,
\[
  r = \left\lceil \frac{16}{\varepsilon} \right\rceil \le \frac{16}{\varepsilon} + 1 \le \frac{32}{\varepsilon}.
\]
Hence
\[
  |S| \ge \frac{k}{r} \ge \frac{n/8}{32/\varepsilon} = \frac{\varepsilon n}{256}.
\]
The construction above was used only under the standing assumptions made at the beginning
(in particular that the graph has at least one edge). If the graph is edgeless, taking $S = V$
is trivially $\varepsilon$-light. It remains to look at small values of $n$: for $1 \le n \le 3$ any
single vertex set $S = \{v\}$ has $L_S = 0$ and hence is $\varepsilon$-light, and it satisfies
$|S| = 1 \ge \varepsilon n/256$ because $\varepsilon \le 1$. The cases $n = 0$ or $\varepsilon = 0$ were disposed of
at the start. Thus in all cases there exists an $\varepsilon$-light set $S$ with
\[
  |S| \ge \frac{\varepsilon}{256}\,|V|.
\]
This proves the statement with the universal constant $c = 1/256$.

\begin{thebibliography}{1}
\bibitem{BSS}
  J.~Batson, D.~A.~Spielman and N.~Srivastava,
  Twice--Ramanujan sparsifiers,
  \textit{SIAM Journal on Computing} \textbf{41} (2012), no.~6, 1704--1721.
\end{thebibliography}

\end{document}
